\section{PCA 从最大方差到最小重构误差}
\label{sec:13-Machine-Learning/04-Dimensionality-Reduction/01}

化工厂的中控室里挂着一块屏，一条产线上 42 个传感器每秒各报一个数：温度、压力、流量、振动。把 42 条曲线叠在同一张图上，除了开车时段一起往上窜，什么也读不出来。运维想要的其实很朴素——用三四条曲线概括这条线今天的状态。于是问题变成：哪三四条？挑波动最大的传感器吗？振动通道数值上动静最大，可它大部分是高频噪声；那支读数沉稳的反应釜温度计，说不定才是产品合格率的关键。

\subsection{降维要回答的不是留几列，而是允许丢掉哪些方向}

先把一个常见的误解摆正。降维不等于``从 42 列里挑 5 列''。挑列是把坐标轴限制在原始变量上，而降维允许你换一组坐标轴——新的方向可以是若干传感器的线性组合，比如``全线整体升温''这个模式，它不对应任何单个探头。

真正需要你签字的是另一件事：\textbf{哪些方向上的信息可以被丢弃}。任何降维方法都必须给出这个声明，区别只在声明的内容。PCA 声明的是：方差小的方向可以丢。随机投影几乎不挑方向，它声明的是点对之间的距离要保住、别的都无所谓。t-SNE 一类的方法声明的是邻居关系要保住，全局距离可以扭曲。把这三条并排看，就明白``降维''从来不是一个动作，而是一族互不相同的取舍。

PCA 的取舍有两个等价的说法。正着说：挑出让样本互相散得最开的正交方向。反着说：挑出一个低维子空间，使得把样本正交投影上去、再投影回来的误差最小。两句话在数学上是同一件事，但它们分别对应两种直觉——一种关心保留了什么，一种关心损失了什么。

\subsection{忘记中心化，第一主成分会指向数据的重心}

把 \(n\) 个样本按行排成 \(X\in\R^{n\times p}\)，一行一个样本，一列一个变量。第一步是去均值：

\[
X_c = X - \ind\,\bar x^{\trans},
\]

其中 \(\ind\) 是全 1 列向量，\(\bar x\) 由各列的样本均值组成。

这一步不是可有可无的清洗动作。PCA 最大化的是 \(\norm{Xv}_2^2\) 这种以\emph{原点}为参照的平方长度。如果数据云整体远离原点，那么``离原点最远''和``彼此散得最开''就是两回事，前者会压倒后者。极端情形下，第一主方向几乎与均值向量 \(\bar x\) 平行：它描述的是这批数据待在哪里，而不是样本之间怎么变化。

\begin{center}
\begin{tikzpicture}[>=stealth,scale=0.92]
% ---- 左：已中心化 ----
\draw[gray!55,->] (-2.9,0) -- (2.9,0);
\draw[gray!55,->] (0,-1.9) -- (0,1.9);
\fill[black] (0,0) circle (1.4pt);
\node[font=\scriptsize,gray!70] at (0.3,-0.28) {原点};
\foreach \x/\y in {-2.0/-0.9, -1.7/-1.1, -1.5/-0.8, -1.2/-0.5, -0.8/-0.6,
                   -0.5/-0.1, -0.2/-0.3, 0.0/-0.1, 0.1/0.2, 0.4/0.0,
                   0.7/0.5, 0.9/0.25, 1.0/0.6, 1.3/0.55, 1.7/0.7, 2.0/1.1}
  \fill[black!62] (\x,\y) circle (1.5pt);
\draw[very thick,->] (-2.3,-1.15) -- (2.3,1.15);
\node[font=\scriptsize] at (2.05,1.62) {\(v_1\)};
\node[font=\small] at (0,-2.5) {中心化之后};

% ---- 右：未中心化 ----
\begin{scope}[shift={(8.4,-1.35)}]
\draw[gray!55,->] (-4.1,0) -- (1.3,0);
\draw[gray!55,->] (0,-0.5) -- (0,4.6);
\fill[black] (0,0) circle (1.4pt);
\node[font=\scriptsize,gray!70] at (0.4,-0.32) {原点};
\foreach \x/\y in {-2.0/-0.9, -1.7/-1.1, -1.5/-0.8, -1.2/-0.5, -0.8/-0.6,
                   -0.5/-0.1, -0.2/-0.3, 0.0/-0.1, 0.1/0.2, 0.4/0.0,
                   0.7/0.5, 0.9/0.25, 1.0/0.6, 1.3/0.55, 1.7/0.7, 2.0/1.1}
  \fill[black!62] ({\x-1.5},{\y+3.0}) circle (1.5pt);
\draw[very thick,->] (0,0) -- (-1.88,3.76);
\node[font=\scriptsize,anchor=east] at (-2.02,3.86) {\(v_1\)};
\draw[dashed,black!55] (-3.2,2.15) -- (0.3,3.90);
\node[font=\scriptsize,black!60,anchor=west] at (0.44,3.78) {真实变化方向};
\fill[black] (-1.5,3.0) circle (1.8pt);
\node[font=\scriptsize] at (-1.15,2.75) {\(\bar x\)};
\node[font=\small] at (-1.3,-1.15) {未中心化};
\end{scope}
\end{tikzpicture}
\end{center}

\emph{同一朵数据云，左边围绕原点，第一主方向沿着云的长轴；右边整体平移之后，第一主方向被拉去指向重心，与真实的变化方向（虚线）几乎垂直。}

图右侧那条粗箭头解释了一个在实践中反复出现的症状。忘记中心化时，你会看到第一成分的解释方差比高得离谱，接近 \(99\%\)；与此同时，所有样本在这个成分上的得分挤在一起，彼此差别小到需要放大坐标轴才看得见。原因就在图里：得分近似等于每个样本在 \(\bar x\) 方向上的投影，而这个投影对所有样本几乎相同，都约等于 \(\norm{\bar x}_2\)。解释方差比在骗你，因为它计算的是相对于原点的平方和，不是相对于均值的平方和。看到这一对症状同时出现，先回头查中心化，不要急着解释 loading。完整的代数推导在线性代数部分的\hyperref[sec:03-Linear-Algebra/07-SVD-and-Data-Applications/04]{``从数据矩阵到 PCA''}一节。

尺度问题紧随其后，而它没有一刀切的答案。协方差 PCA 用的是中心化后的原始数值，即 \(\frac{1}{n-1}X_c^{\trans}X_c\) 的谱。身高以米计、收入以元计时，协方差阵几乎完全被收入统治，第一主成分不过是``收入方向''的换个说法。标准化到单位方差——等价于对相关矩阵做 PCA——让每个变量拿到同等的发言权，同时也把原本方差很小的噪声变量一并放大。变量量纲一致且方差量级本身有科学含义时，标准化会稀释真实的结构；量纲各异时，不标准化会让大数压死小数。这是一个建模决定，不是一个默认参数。

\subsection{最大方差与最小重构误差是同一句话的两种说法}

沿单位方向 \(v\) 的投影得分是 \(X_c v\)，其未归一化的样本方差为 \(v^{\trans}X_c^{\trans}X_c v\)。第一主方向因此是

\[
\max_{\norm{v}_2=1}\ v^{\trans} X_c^{\trans} X_c v
\]

的解。这是 Rayleigh 商的标准形式，最优 \(v_1\) 取 \(X_c^{\trans}X_c\) 最大特征值对应的特征向量；后续方向在与已选方向正交的约束下依次最大化剩余方差，这一串结论都由实对称矩阵的谱定理承担。

计算上更稳妥的路线是直接对 \(X_c\) 做 SVD：写 \(X_c = U\Sigma V^{\trans}\)，则 \(V\) 的列是主方向（loadings），\(X_cV = U\Sigma\) 是样本得分（scores），第 \(j\) 个方向解释的方差为 \(\sigma_j^2/(n-1)\)。绕道 \(X_c^{\trans}X_c\) 会把条件数平方，下一节会说清这件事。

保留前 \(k\) 个方向组成 \(V_k\)，低维表示是 \(Z = X_cV_k\)，重构回原空间得到 \(\widehat X_c = ZV_k^{\trans} = X_cV_kV_k^{\trans}\)。这里 \(V_kV_k^{\trans}\) 是到主子空间的正交投影矩阵，于是总平方长度可以按勾股关系拆开：

\[
\begin{aligned}
\norm{X_c}_F^2
&= \norm{X_cV_kV_k^{\trans}}_F^2 + \norm{X_c - X_cV_kV_k^{\trans}}_F^2\\
&= \sum_{j\le k}\sigma_j^2 + \sum_{j>k}\sigma_j^2 .
\end{aligned}
\]

左端与 \(k\) 无关。总量固定，多留一点就少丢一点，两个目标因此严格等价，而丢掉的部分恰好是被截断的奇异值平方和。这也说明``解释方差比''本质上是奇异值平方的份额，不是别的什么神秘量。

等价性同时画出了 PCA 的边界。它最优的是\emph{线性子空间下的平方重构}，仅此而已。方差小的方向可能恰好携带疾病分型的信息；方差大的方向可能只是光照变化或批次效应。构造一个反例并不困难：两类二维数据，共同漂移的方向方差极大，区分两类的方向方差很小；第一主成分几乎与漂移方向重合，只留它做分类，两类的投影严重重叠，还不如直接用原始两维。无监督重构的最优方向，和监督预测的有用方向，是两个不同的最优。

\subsection{碎石图上的肘部可能完全是噪声画出来的}

拿到奇异值序列之后，通常的做法是画碎石图找肘部。这个做法在 \(n\) 远大于 \(p\) 时还算可靠，在高维下则要小心：即使数据完全是各向同性的白噪声，样本协方差的特征值也不会挤在一起，而是散布成一条有宽度的分布，最大的那个明显高于总体值。谱的失真是系统性的，不是运气不好。于是碎石图上出现一段下降、然后趋平的形状，肘部看起来煞有介事，实际可能纯由采样涨落造出来。

判断信号的正确基准不是``看起来像肘部''，而是随机矩阵理论给出的零模型上界。当维数比 \(p/n\to\gamma\) 时，纯噪声情形下样本特征值几乎必然落在 Marchenko--Pastur 分布的支撑内，其上界为

\[
\lambda_+ = \sigma^2\bigl(1+\sqrt{\gamma}\bigr)^2 .
\]

只有明显超出 \(\lambda_+\) 的特征值，才有资格被当作信号；落在支撑内的那些，无论碎石图上看着多像一道台阶，都不构成证据。这条基准以及它在有真实信号时的位移规律，见\hyperref[ch:094]{``随机矩阵理论：高维协方差的谱失真''}一章，尤其是\hyperref[sec:11-Mathematical-Statistics/09-Random-Matrix-Theory/02]{``Marchenko--Pastur 律给出谱的精确形状''}一节。

选 \(k\) 时还有几件事值得一并看：谱隙是否明显、保留的成分能否满足压缩率或去噪质量的要求、下游验证指标在哪个 \(k\) 之后停止改善、主子空间在子抽样下是否稳定。最后一条尤其重要。单个 loading 的符号没有固有含义，\(v\) 与 \(-v\) 是同一条轴；奇异值接近时，那个子空间内部的具体基会随数据扰动自由旋转，此时应当解释整个子空间，而不是纠缠某一列的数值。

\subsection{降维之后再训练，你已经换了一个问题}

把 PCA 接到预测流程前面，看上去只是加了一步预处理，实际上改变了问题本身。模型不再在原始变量上学习，而是在一组由无监督准则挑出来的坐标上学习；这组坐标的选择完全没有参考标签。如果标签信息恰好藏在被丢弃的小方差方向里，任何下游模型都无法把它找回来——信息在进入模型之前就没了。降维带来的收益（更少的参数、更低的方差、更快的训练）和这份不可逆的损失，需要放在一起衡量。

由此引出一条纪律：均值、标准差、主方向，全部只能从训练折估计，然后原样应用到验证折与测试集。在全数据上先跑一遍 PCA 再划分折，验证样本的结构已经参与了主方向的估计，得到的分数会偏乐观，而且这种乐观很隐蔽——它不像标签泄漏那样明显，却同样会让线上表现低于离线评估。判断标准始终是：交叉验证评估的对象是整条训练流程，不是流程末端的那个模型，见\hyperref[sec:11-Mathematical-Statistics/06-Resampling-and-Model-Selection/03]{``交叉验证评估的是完整训练流程''}一节。

到这里，PCA 的最优性还只是在``所有正交投影''这个范围内说的。下一节把范围放大到所有低秩矩阵，Eckart--Young--Mirsky 定理会说明截断依然最优——同时说明这份最优在什么条件下立刻失效。
